\documentclass[a4paper,twoside]{article}

\usepackage{morewrites}

\usepackage[svgnames,dvipsnames,x11names]{xcolor}
\usepackage{graphicx}
\usepackage[english]{babel}
\usepackage[colorlinks=true, linkcolor=NavyBlue, urlcolor=NavyBlue, citecolor=NavyBlue]{hyperref}
\usepackage[a4paper]{geometry}
\usepackage{ts-fonts}
\usepackage{tcolorbox}
\usepackage{fvextra}
\usepackage{amsmath}
\usepackage{float}
\usepackage{seqsplit}
\usepackage{ts-code}

\usepackage{lastpage}
\usepackage{refcount}
\makeatletter
\def\ps@plain{
\let\@oddhead\@empty
\let\@evenhead\@empty
\def\@oddfoot{
\hfil
\ifnum\getpagerefnumber{LastPage}>1\relax
\thepage
\fi
\hfil}
\let\@evenfoot\@oddfoot
}
\makeatother

\title{Wiki Links}
\author{}\date{}
\begin{document}
\maketitle
MkDocs and Python-Markdown support the \tscodeinline{[[Wiki Link]]} syntax via the \tscodeinline{wikilinks}
extension. TeXSmith keeps that behavior so you can link between pages without
remembering exact file paths.
\begin{tscode}[lang=markdown, engine=pygments]
\begin{Verbatim}[commandchars=\\\{\},breaklines, breakanywhere, commandchars=\\\{\}]
[[Getting Started]]
[[Subfolder/Page Title|Custom label]]
\end{Verbatim}
\end{tscode}
\begin{itemize}
\item The portion before the pipe resolves to a Markdown file (\tscodeinline{Getting Started} \tsscript{symbols}{\(\rightarrow\)}
\tscodeinline{getting-\allowbreak{}started.md}).
\item Anything after \tscodeinline{|} becomes the rendered link text.
\item When building PDFs, TeXSmith turns wiki links into standard hyperlinks, so the
references remain navigable.
\end{itemize}
\end{document}
